\documentclass[conference]{IEEEtran}
\IEEEoverridecommandlockouts
% The preceding line is only needed to identify funding in the first footnote. If that is unneeded, please comment it out.
\usepackage{cite}
\usepackage{amsmath,amssymb,amsfonts}
\usepackage{algorithmic}
\usepackage{graphicx}
\usepackage{textcomp}
\usepackage{xcolor}
\usepackage{hyperref}
\def\BibTeX{{\rm B\kern-.05em{\sc i\kern-.025em b}\kern-.08em
    T\kern-.1667em\lower.7ex\hbox{E}\kern-.125emX}}
\begin{document}

\title{Multi-Robot Coverage Path Planning using DARP on Crazyflie Drones}

\author{\IEEEauthorblockN{1\textsuperscript{st} Jorge Prieto Álvarez}
\IEEEauthorblockA{\textit{MSc. in Electronics, Robotics and Automation Engineering} \\
\textit{Universidad de Sevilla}\\
Seville, Spain}
\and
\IEEEauthorblockN{2\textsuperscript{nd} Víctor Javier Granero Gil}
\IEEEauthorblockA{\textit{MSc. in Electronics, Robotics and Automation Engineering} \\
\textit{Universidad de Sevilla}\\
Seville, Spain}
\and
\IEEEauthorblockN{3\textsuperscript{rd} José Francisco López Ruiz}
\IEEEauthorblockA{\textit{MSc. in Electronics, Robotics and Automation Engineering} \\
\textit{Universidad de Sevilla}\\
Seville, Spain}
}

\maketitle

\begin{abstract}
This paper presents a comprehensive implementation of a Multi-Robot Coverage Path Planning (CPP) system using the Divide Areas Algorithm for Robotics Platforms (DARP). The system is designed for Crazyflie 2.1 drones and is implemented using the Robot Operating System 2 (ROS2). The architecture leverages a modular approach with dedicated nodes for path calculation, mission planning, and flight control. The proposed solution is validated through simulations in Gazebo, demonstrating efficient area coverage and obstacle avoidance capabilities in a multi-agent scenario.
\end{abstract}

\begin{IEEEkeywords}
Coverage Path Planning, DARP, ROS2, Crazyflie, Multi-Robot Systems, Gazebo
\end{IEEEkeywords}

\section{Introduction}
Coverage Path Planning (CPP) is a fundamental problem in mobile robotics, where the goal is to determine a path that passes over all points of an area of interest while avoiding obstacles. This capability is essential for various applications, including search and rescue operations, precision agriculture, floor cleaning, and environmental monitoring.

Multi-robot systems offer significant advantages over single-robot approaches, such as reduced execution time and increased robustness through redundancy. However, coordinating multiple robots introduces challenges in workload distribution and collision avoidance.

This project addresses these challenges by integrating the Divide Areas Algorithm for Robotics Platforms (DARP) \cite{article} with the Crazyflie 2.1 drone platform. We utilize ROS2 \cite{macenski2022ros2} for inter-process communication and Gazebo \cite{koenig2004gazebo} for high-fidelity simulation. The system is designed to be modular, allowing for easy scalability and adaptation to different environments.

The remainder of this paper is organized as follows: Section II reviews the state of the art in CPP. Section III details the system implementation and architecture. Section IV presents the simulation setup and results. Finally, Section V concludes the paper and discusses future work.

\section{State of the Art}
The CPP problem has been extensively studied in the literature. Approaches can be broadly categorized into heuristic, grid-based, and cellular decomposition methods.

Grid-based methods decompose the environment into a collection of uniform cells. This simplifies the representation but can lead to exponential complexity in high-resolution maps. Spanning Tree Coverage (STC) is a popular grid-based algorithm that guarantees complete coverage but is typically limited to single-robot scenarios or requires complex modifications for multi-robot teams.

For multi-robot systems, the primary challenge is partitioning the area efficiently. The DARP algorithm, proposed by Kapoutsis et al., addresses this by dividing the grid into equal (or weighted) portions for each robot based on their initial positions. It uses a cyclic gradient optimization to ensure that each robot's assigned area is connected and obstacle-free. Once the area is partitioned, a spanning tree algorithm is used to generate the coverage path for each robot within its assigned zone.

This project adopts DARP due to its ability to guarantee complete coverage, non-overlapping areas, and connectivity, which are critical for efficient multi-UAV operations.

\section{Implementation}

\subsection{System Architecture}
The system is built upon ROS2 (Robot Operating System 2) and utilizes the \texttt{crazyswarm2} architecture for interfacing with the Crazyflie drones. The core logic is distributed across three custom ROS2 nodes: \texttt{darp\_node}, \texttt{planning\_node}, and \texttt{control\_node}. Figure \ref{fig:arch} (conceptual) illustrates the data flow between these components.

\subsection{DARP Node}
The \texttt{darp\_node} encapsulates the DARP algorithm logic. It exposes a service \texttt{darp\_service} that accepts the mission parameters:
\begin{itemize}
    \item Operational area dimensions (min/max X and Y).
    \item Obstacle positions.
    \item Initial positions of the robots.
\end{itemize}

Upon receiving a request, the node discretizes the continuous world coordinates into a grid based on a configurable \texttt{cell\_size}. It then executes the DARP algorithm to partition the grid and generate paths. The output consists of:
\begin{itemize}
    \item \textbf{Trajectories}: A list of waypoints for each robot.
    \item \textbf{Zones}: A grid map indicating the assigned area for each robot, useful for visualization.
\end{itemize}

The node handles the coordinate transformation between the metric world frame and the discrete grid indices required by DARP.

\subsection{Planning Node}
The \texttt{planning\_node} acts as the central coordinator. Its responsibilities include:
\begin{itemize}
    \item \textbf{State Management}: It monitors the positions of all agents via \texttt{/agent\_id/state/position} topics.
    \item \textbf{Mission Triggering}: It waits for all agents to be ready before requesting a plan from the \texttt{darp\_node}.
    \item \textbf{Trajectory Distribution}: Once a solution is received, it publishes specific trajectories to each agent on \texttt{/agent\_id/planning/trajectory}.
    \item \textbf{Dynamic Replanning}: The node supports replanning capabilities. It tracks visited cells and can treat them as obstacles in subsequent DARP requests, ensuring that previously covered areas are not repeated if a mission update is required.
\end{itemize}

\subsection{Control Node}
The \texttt{control\_node} runs an instance for each robot and manages the flight state machine. It subscribes to the assigned trajectory and the robot's current position. The internal state machine consists of the following states:
\begin{enumerate}
    \item \textbf{INIT}: Sends initial setpoints to prepare for takeoff.
    \item \textbf{TAKEOFF}: Executes the takeoff maneuver to a target altitude.
    \item \textbf{MISSION}: Follows the received trajectory waypoints. It implements a "stop-and-go" or continuous path following logic, switching to the next waypoint when the robot is within an \texttt{acceptance\_radius} of the current target.
    \item \textbf{LANDING}: Initiates landing once the trajectory is completed.
    \item \textbf{LANDED}: Disarms the drone and terminates the mission.
\end{enumerate}

This node interfaces with the lower-level Crazyflie controllers via standard ROS2 topics and services (Takeoff, Land, Arm).

\section{Simulation and Results}

\subsection{Simulation Environment}
The system was validated using the Gazebo simulator integrated with \texttt{crazyswarm2}. The simulation environment was configured with:
\begin{itemize}
    \item A defined workspace (e.g., $10 \times 10$ meters).
    \item Multiple Crazyflie 2.1 models.
    \item Static obstacles placed within the area.
\end{itemize}

\subsection{Results}
The experiments demonstrated the effectiveness of the DARP integration. 
\begin{itemize}
    \item \textbf{Area Partitioning}: The algorithm successfully divided the available free space into distinct regions for each drone.
    \item \textbf{Coverage}: The generated paths covered all accessible cells within the assigned zones.
    \item \textbf{Execution}: The \texttt{control\_node} successfully guided the drones through the generated waypoints with stable flight behavior.
    \item \textbf{Replanning}: The system showed the capability to update plans dynamically by marking visited cells, although the primary focus was on static offline planning.
\end{itemize}

Visualizations in RViZ confirmed the alignment of the planned grid zones with the real-world positions of the drones.

\section{Conclusion}
This paper presented a ROS2-based framework for multi-robot coverage path planning using Crazyflie drones. By integrating the DARP algorithm, we achieved an efficient division of labor among multiple agents. The modular architecture allows for easy extension, such as adding dynamic obstacle avoidance or scaling to a larger swarm. Future work will focus on deploying this system on physical Crazyflie hardware and implementing real-time decentralized coordination.

\section*{Acknowledgment}
The authors would like to thank the teaching staff of the Robotics Projects course for their guidance and support.

\bibliographystyle{IEEEtran}
\bibliography{bibliography}

\end{document}